\documentclass[10pt,a4paper]{amsart}
\usepackage[margin=19mm]{geometry}
\usepackage{amsmath,amssymb,mathtools}
\usepackage[scr=rsfs,cal=euler]{mathalfa}
\usepackage{xfrac}
\usepackage[unicode,hidelinks,bookmarks=false]{hyperref}
\newcommand{\ie}{i.e.\ }
\newcommand{\cf}{cf.\ }
\newcommand{\resp}{respectively\ }
\newcommand{\Subsection}{\subsection}
\providecommand{\nobreakdash}{\nobreak\hbox{-}}
\setlength{\emergencystretch}{2em}
\multlinegap0pt
\usepackage{amssymb}
\usepackage{mathtools}
\usepackage{bm}
\usepackage[shortlabels]{enumitem}
\usepackage{xcolor}
\usepackage{algorithm}\usepackage{algpseudocode}
\newtheorem{theorem}{Theorem}
\title{PAGR: Proof-Carrying Algebraic-Geometric Retrieval:\\ A Quiver-, Provenance-, and Sheaf-Theoretic Framework for Grounded LLM Retrieval}
\author{Xingting Wang, Min Wu}
\date{Source: arXiv:2609.06127v1 (CC BY 4.0)}
\begin{document}
\maketitle

\noindent\emph{Source and licence.} PAGR: Proof-Carrying Algebraic-Geometric Retrieval:\\ A Quiver-, Provenance-, and Sheaf-Theoretic Framework for Grounded LLM Retrieval. Version arXiv:2609.06127v1 (CC BY 4.0), distributed by arXiv under the Creative Commons Attribution 4.0 International licence, http://creativecommons.org/licenses/by/4.0/, as recorded in the arXiv licence metadata for this exact version and shown on the abstract page https://arxiv.org/abs/2609.06127v1. Full licence notice: \texttt{licenses/arxiv-2609.06127-CC-BY-4.0.txt}.\par\smallskip

\noindent\textit{Editorial scope.} Counted unit: the article's theorem thm:margin (source0.tex lines 513--524). The article's complete original proof of it is delivered with it (source0.tex lines 525--533, one \texttt{proof} environment of 9 lines). No mathematics is written, repaired or re-derived: the statement and the proof are the article's own lines, in the article's own notation, and no retained line is substituted or re-flowed.  No further source-local context is needed: every cross-reference of every retained line resolves inside the two delivered spans.  Nothing was omitted from any retained span, so the package carries no omission record.  No retained line contains a citation, so the wrapper carries no bibliography and no bibliography entry.  No retained line contains a figure environment, an image-inclusion command or drawing code, so no image is delivered and the package declares no asset dependency.  Editorial formatting only, in addition: an AMS \texttt{amsart} preamble that restates the article's own declarations and macros used by the retained lines, namely the environments \texttt{theorem} on separate counters, together with the standard packages those lines need.  The retained environments print in this document as Theorem 1 in order of appearance, whereas the article prints the same items with the numbering of its own full document.  The complete native source of the article as distributed in the arXiv e-print for this exact version is bundled unchanged as \texttt{sources/209492/source0.tex}.
\begin{theorem}[Margin stability of geometric seeding]
\label{thm:margin}
Assume a perturbation changes every entity score by at most $\varepsilon$:
\[
|s_q'(x)-s_q(x)|\le \varepsilon.
\]
Let $x_{(k)}$ and $x_{(k+1)}$ denote the $k$th and $(k+1)$st ranked entities under $s_q$. If
\[
s_q(x_{(k)})-s_q(x_{(k+1)})>2\varepsilon,
\]
then the top-$k$ seed set is unchanged under the perturbation.
\end{theorem}

\begin{proof}
The strict margin implies that the top-$k$ set is uniquely determined under $s_q$. Let $x$ be ranked in the top $k$ and $x'$ outside it, so that $s_q(x)\ge s_q(x_{(k)})$ and $s_q(x')\le s_q(x_{(k+1)})$. Then
\[
s_q'(x)\ \ge\ s_q(x)-\varepsilon\ \ge\ s_q(x_{(k)})-\varepsilon
\qquad\text{and}\qquad
s_q'(x')\ \le\ s_q(x')+\varepsilon\ \le\ s_q(x_{(k+1)})+\varepsilon .
\]
The hypothesis $s_q(x_{(k)})-s_q(x_{(k+1)})>2\varepsilon$ rearranges to $s_q(x_{(k)})-\varepsilon>s_q(x_{(k+1)})+\varepsilon$, hence $s_q'(x)>s_q'(x')$ for every such pair. No boundary crossing occurs and the top-$k$ set is unchanged.
\end{proof}

\end{document}
